\section{Scope, Method, and Qualifications}

\subsection*{Audience and Purpose}

Intelligent adult readers---with no prerequisite training in scholastic systematics or symbolic logic---who approach classical Catholic metaphysics sympathetically yet refuse to substitute deference for understanding are the intended audience. The sole objective is practical competence: enabling a reader to open an unmarked page of Aristotle, St.\ Thomas Aquinas, or a serious philosophical interlocutor with a pencil and determine what is being claimed, what the load-bearing terms mean, what warrant supports each premise, where the substantive philosophical burden lies, what follows and what does not yet follow, and what rejecting the burden costs.

What follows is not a comprehensive system of Thomistic philosophy, a manual of formal logic, a history of medieval schools, or an apologetic proving that every classical argument succeeds. Where an examination reveals that a passage proves less than readers commonly assume, or rests on a contested metaphysical principle whose denial constitutes a coherent rival framework rather than a logical contradiction, that finding is reported plainly as an achievement of honest reading.

\subsection*{Methodological Boundaries and Historical Discipline}

The six-question reading frame (\textbf{Claim, Meaning, Warrant, Burden, Consequence, Denial}), the working list of eight warrant kinds, and the five kinds of denial are modern pedagogical instruments designed to make classical distinctions visible to readers without scholastic formation. They are never attributed to Aristotle, Aquinas, or their commentators. Where classical theory is presented, it is presented in its own historical terms (Aristotle's six conditions on demonstrative premises, \emph{episteme}, \emph{apodeixis}, \emph{endoxa}, and \emph{quia}/\emph{propter quid}; Aquinas's \emph{duplex est demonstratio}, \emph{per se notum quoad se}/\emph{quoad nos}, \emph{praeambula fidei}, \emph{res significata}/\emph{modus significandi}, \emph{separatio}, and \emph{per se}/\emph{per accidens} causal ordering).

Prose is treated as canonical throughout. Displayed symbolic notation appears strictly where it isolates a named scope shift, modal ambiguity, identity distinction, formal invalidity, or competing formal reading, accompanied in every case by an explicit licensing sentence. Metaphysical notions such as act, potency, essence, existence, analogy, and causation are never routinely encoded as predicate letters.

\subsection*{Faith, Natural Reason, and Ecclesial Boundaries}

Within Catholic theology, revealed articles of faith and conclusions of natural reason are distinct without being opposed. Throughout the nine lessons, premises held on divine revelation (such as the temporal beginning of the world in \emph{ST} I, q.\,46, a.\,2) remain visibly labelled as revelation: they are neither converted into demonstrations of natural reason nor faulted for lacking proofs they never claimed. Conversely, philosophical claims offered as accessible to natural reason are evaluated strictly by philosophical warrant and never defended by appeal to saintly or authorial prestige. Offered strictly as a personal study aid in philosophical reading, the material carries no magisterial authority, does not adjudicate the theology of revelation, and stands as neither an official textbook nor a critical edition.

\subsection*{Editions, Translations, and Rights Boundaries}

Quotations from St.\ Thomas Aquinas's \emph{Summa Theologiae} (Prima Pars) are checked against the Leonine Latin text (Vol.\ IV, Rome, 1888, via Corpus Thomisticum) and quoted in English from the public-domain translation of the Fathers of the English Dominican Province (Benziger Brothers; Project Gutenberg eBook 17611), except where explicitly marked as a close rendering preserving Latin clause structure. Aquinas's \emph{De ente et essentia} is checked against the Latin edition of L.\ Baur emended by J.\ Koch (M\"unster, 1933, via Corpus Thomisticum). Aristotle's \emph{Posterior Analytics}, \emph{Prior Analytics}, \emph{Metaphysics}, \emph{Nicomachean Ethics}, and \emph{Topics} are checked against the public-domain Oxford translations of G.\ R.\ G.\ Mure, A.\ J.\ Jenkinson, W.\ D.\ Ross, and W.\ A.\ Pickard-Cambridge, governed by Bekker pagination. Bl.\ John Duns Scotus's \emph{Ordinatio} I, d.\,3, p.\,1, q.\,1--2, n.\,26 is checked against the Vatican edition text; David Hume's \emph{Dialogues Concerning Natural Religion} (1779, Part IX) is quoted from the public-domain text. Unverified loci---including \emph{In Phys.}\ VIII, \emph{SCG} II, c.\,38, and \emph{De ver.}\ q.\,10, a.\,12---are nowhere cited as support for claims here.

Public-domain primary texts and translations remain in the public domain; third-party scholarly citations carry no transfer of rights. Original prose, pedagogy, and analysis composed here fall under the CC BY 4.0 license.

\subsection*{Review Boundaries}

Primary quotations from Aquinas's \emph{Summa Theologiae} (Prima Pars) and \emph{De ente et essentia} were verified against tracked local witnesses, and the six-question reading method was tested internally on unseen articles of Aquinas (\emph{ST} I, q.\,8, a.\,1 and q.\,9, a.\,1). No external specialist, clerical, or peer review is claimed.
